% !TEX root = ../main.tex
\section{Noah}
Empirical evaluations in evolutionary computing require rigorous statistical reporting and transparent procedural documentation.
Sequential experimental protocols should be structured using numbered enumeration to communicate order of execution clearly.

\subsection{Numbered Sequences and Procedures}
When documenting iterative experimental benchmarks, numbered lists establish the exact procedural sequence:
\begin{enumerate}
    \item Random seeds and population vectors were initialized uniformly across sample spaces.
    \item Fitness evaluations ran over 50 independent runs for each algorithmic variant.
    \item Non-parametric statistical tests were executed to evaluate convergence speed.
\end{enumerate}

\subsection{Statistical and Mathematical Notation}
APA 7 mandates specific typographic rules for statistical reporting.
All single-letter statistical symbols and variables—such as \(p\), \(t\), \(F\), \(M\), \(SD\), and \(d\)—must appear in italics.
Numbers, parentheses, and Greek letters like \(\alpha\) or \(\eta_p^2\) remain upright.
Furthermore, numbers bounded between 0 and 1 that cannot theoretically exceed 1.0 (including \(p\)-values, correlation coefficients \(r\), and proportions) omit the leading zero.

The evolutionary strategy demonstrated superior fitness over the baseline, \(M = 94.20\), \(SD = 3.15\), compared to the control group, \(M = 88.50\), \(SD = 4.22\).
An independent samples \(t\)-test revealed a statistically significant difference, \(t(48) = 5.41\), \(p < .001\), \(d = 1.53\), with a 95\% confidence interval around the difference, \(95\% \text{ CI } [3.58, 7.82]\).
Similarly, analysis of variance confirmed main effects across benchmark landscapes, \(F(2, 72) = 8.64\), \(p = .002\), \(\eta_p^2 = .19\).
When accompanied by specific numerals, format percentages using the symbol, such as 95\% confidence and 12\% variance, but use the word \enquote{percentage} in abstract text.
